\section{Advanced tool-wide runtime options}
\label{runtime}

The command-line tools support a set of advanced options that apply uniformly
across all tools. These options control tool-wide runtime behavior and are
intentionally documented separately from individual manual pages, since they
are not specific to any single utility.

The options described in this section adjust internal runtime settings for which
reasonable defaults are provided. Most users will never need to change these
defaults, but advanced users may wish to tune them to better match available
system resources or specific processing workloads.

All options are specified using a {\file -J} prefix, which causes them to be
passed directly to the Java virtual machine rather than processed as
tool-specific command-line arguments. This mechanism allows the options to
modify the Java execution environment in which the tool runs, and therefore
affect behavior consistently across the entire tool suite.

The available options are described below:

\begin{description}[style=nextline]

\item[{\file -J-XmxSIZE}] 

Controls the maximum heap memory available to the Java virtual machine. This
setting places an upper limit on the amount of memory the tool may use at
runtime.

Reasonable defaults are selected automatically based on the tool being
executed, and most users do not need to specify this option. Tools with
more memory-intensive workloads use larger default heap sizes; for example,
{\file cwrender} allocates more heap memory by default than {\file cwinfo}.
Increasing the heap size may be beneficial when processing very large datasets,
building large composites, or performing memory-intensive operations.

Example: {\file -J-Xmx8g}

\item[{\file -J-Dcw.cache.size=SIZE}] 

Controls the maximum size of the in-memory cache used to store tiles of
two-dimensional data read from input files. This option is currently implemented
for NetCDF input files only.

The cache size is specified in megabytes, with a default value of 512~Mb.
Increasing the cache size may be useful when working with files written using
monolithic compression, where entire variables are compressed as a single chunk.
In such cases, even small data accesses may require decompressing the full
variable into memory, increasing memory usage and reducing performance during
interactive display or automated processing.

Example: {\file -J-Dcw.cache.size=1024}

\item[{\file -J-Dcw.compress.mode=true$|$false}] 

Controls whether compression is used when writing CoastWatch HDF data files.
The default value is {\file true}.

When enabled, HDF4 data variables are written using square chunked compression,
which improves performance for interactive access and partial reads. In some
workflows, compression may be unnecessary or undesirable. For example, when a
file is written as an intermediate result in a longer computation and will be
deleted afterward, disabling compression can reduce write time and improve
overall processing speed.

When set to {\file false}, HDF4 data files are written without chunking or
compression.

Example: {\file -J-Dcw.compress.mode=false}

\item[{\file -J-Dcw.chunk.size=SIZE}] 

Controls the chunk size, in kilobytes, used when writing CoastWatch HDF data
files. The default value is 512~kb.

This setting specifies the total byte size of each chunk, not the chunk
dimensions. Chunk dimensions are determined automatically based on the data
type and selected chunk size. For example, with a 512~kb chunk size, 16-bit
integer data is written in 512$\times$512 chunks, since
2~bytes $\times$ 512 $\times$ 512 $=$ 524{,}288~bytes. This choice reflects the
historical characteristics of CoastWatch AVHRR data and provides a balance
between efficient I/O and interactive performance.

Using the same chunk size, 8-bit byte data is written in
724$\times$724 chunks, 32-bit floating-point data in 362$\times$362 chunks, and
so on.

Example: {\file -J-Dcw.chunk.size=1024}

\item[{\file -J-Dcw.debug=LEVEL}] 

Controls the verbosity of diagnostic logging output. Valid values
include {\file fine} and {\file finer}, with higher levels producing more
detailed messages.

By default, tools emit only warning and error messages. Enabling debug logging
is useful when diagnosing unexpected behavior, performance issues, or data
handling problems.

This option affects logging output only and does not change processing results.

\item[{\file -J-Dcw.log.format=FORMAT}] 

Controls the format used for log messages written to the console. The default
format is compact and intended for routine use.

Setting this option to {\file long} enables more detailed log output, including
additional contextual information useful during debugging or development.

\item[{\file -J-Dcw.memory.monitor=true}] 

Enables periodic reporting of memory usage during tool execution. When enabled,
memory statistics are printed at intervals to help track resource consumption.

This option is primarily intended for performance analysis and troubleshooting.
It has no effect on processing results.

\item[{\file-J-Dsun.java2d.uiScale=SCALE}] 

Controls the scaling factor used for graphical user interface elements in
Java-based graphical tools.

This option may be useful on high-DPI displays or systems with display scaling
enabled, where default UI sizing results in text or controls appearing too
small.

Typical values range from {\file 1.0} to {\file 2.0}.

\end{description}


